\documentclass[12pt]{article}
\usepackage[utf8]{inputenc}
\usepackage[T1]{fontenc}
\usepackage{hyperref}
\title{Mundari Language Resources for K-12 Curriculum Development}
\author{Compiled by Claude Sonnet 5 (claude-sonnet-5) for Andr\'as Kornai}
\date{}
\begin{document}
\maketitle
\section{Introduction}
Mundari (ISO 639-3 unr) is a Munda (Austroasiatic) language of eastern India, closely related to Santali. Take care not to confuse it with the unrelated Mundari of South Sudan, whose Latin-script dictionary also turns up in searches, or with Mundani of Cameroon. Neither is listed below.
\section{Primary Educational Resources}
The CIIL-NCERT \emph{Mundari Primer} (2024) is a basal reader of the alphabet and numerals for Balvatika/Anganwadi children and adult literacy. \url{https://www.ciil.gov.in/primers/Mundari_Primer.pdf}
\section{Linguistic and Reference Materials}
A Unicode proposal for Mundari Bani describes the script, invented by Rohidas Singh Nag. It reports that about 10\% of speakers can read it and that schools and workshops teach it in Odisha, West Bengal and Jharkhand. \url{https://www.unicode.org/L2/L2021/21031r-mundari.pdf} A short overview is at \url{https://www.endangeredalphabets.net/?p=4431}. Wikipedia mentions older works such as Hoffmann's 1903 grammar. \url{https://wikipedia.com/wiki/ISO_639:unr}
\section{Religious and Community Resources}
None were found in these searches.
\section{Digital and Audio Resources}
None were found in these searches. The Unicode encoding of Mundari Bani is proposed but I have not confirmed its status.
\section{Curriculum Development Context}
The CIIL primer is a usable early-grade starting point. Script choice (Mundari Bani, Devanagari, Latin or others) needs to be settled locally. I did not find a dictionary or textbooks for later grades.
\end{document}